% Resume: Arumalla Mohit Krishna  
% Compile with pdfLaTeX.  

\documentclass[letterpaper,10pt]{article}  

\usepackage[margin=0.5in]{geometry}  
\usepackage[T1]{fontenc}  
\usepackage{lmodern}  
\usepackage{microtype}  
\usepackage{titlesec}  
\usepackage{enumitem}  
\usepackage[hidelinks]{hyperref}  
\usepackage{tabularx}  
\input{glyphtounicode}  

\pdfgentounicode=1  
\pagestyle{empty}  
\urlstyle{same}  
\raggedright  
\setlength{\parindent}{0pt}  
\setlength{\tabcolsep}{0pt}  

\titleformat{\section}  
  {\large\scshape}{}{0em}{}[\titlerule]  

\titlespacing*{\section}{0pt}{6pt}{3pt}  

\newcommand{\entry}[2]{%  
  \noindent  
  \begin{tabularx}{\linewidth}  
    {@{}>{\raggedright\arraybackslash}X@{\hspace{6pt}}r@{}}  
    \textbf{#1} & #2  
  \end{tabularx}\par  
}  

\newcommand{\bullets}{  
  \begin{itemize}[  
    leftmargin=1.2em,  
    topsep=2pt,  
    itemsep=1pt,  
    parsep=0pt,  
    partopsep=0pt  
  ]  
}  

\newcommand{\finishbullets}{  
  \end{itemize}  
}  

\begin{document}  

%----------HEADING----------  
\begin{center}  
    {\LARGE\bfseries Arumalla Mohit Krishna}\\[3pt]  
    +1 (631) 352-8643 $|$  
    \href{mailto:mohitkrishna.arumalla@stonybrook.edu}  
    {mohitkrishna.arumalla@stonybrook.edu} $|$  
    \href{https://linkedin.com/in/mohit-krishna-arumalla}  
    {LinkedIn}  
\end{center}  
\vspace{-8pt}  

%-----------EDUCATION-----------  
\section{Education}  

\entry  
  {\textbf{Stony Brook University}}  
  {\textbf{Aug 2025 -- Dec 2026}}  

M.S. in Computer Science \hfill New York, USA\\  
\textit{Coursework: Natural Language Processing, Introduction to Computer Vision, Vision Language Modeling}  
\vspace{3pt}  

\entry  
  {\textbf{Birla Institute of Technology and Science, Pilani}}  
  {\textbf{Aug 2019 -- 2023}}  

B.E. in Electronics \& Communication; GPA: 8.02/10 \hfill India  

%-----------PROFESSIONAL SUMMARY-----------  
\section{Professional Summary}  

\textbf{Entry-level software engineer and Java full stack developer} with 2+ years of experience building distributed services, semantic search, and real-time data pipelines, actively pursuing an M.S. in Computer Science. Skilled in \textbf{Java, Python, cloud infrastructure, and modern development frameworks}. Proven ability to deliver \textbf{production-ready code} with expertise in \textbf{microservices architecture, API development}, and hands-on experience with \textbf{data engineering pipelines} and \textbf{containerized deployments}. Actively committed to bridging the gap between academic knowledge and industry-ready skills.  

%-----------TECHNICAL SKILLS-----------  
\section{Technical Skills}  

\textbf{Languages:}  
Java, Python, Go, C++, SQL \\  

\textbf{Frameworks \& Tools:}  
Spring Boot, gRPC, Kafka, Apache Flink, Airflow, Docker, Kubernetes, Istio/Envoy \\  

\textbf{Databases \& Search:}  
PostgreSQL, Redis, Elasticsearch, BigQuery \\  

\textbf{Machine Learning \& Data Science:}  
PyTorch, TensorFlow, scikit-learn, Transformers, Vision Transformers, Self-Supervised Learning \\  

\textbf{Cloud \& Infrastructure:}  
AWS, NVIDIA Triton, CI/CD pipelines, Linux, Git \\  

%-----------EXPERIENCE-----------  
\section{Experience}  

\entry  
  {\textbf{Software Engineer II} $|$ \textbf{Infra.Market, Bangalore}}  
  {\textbf{Apr 2024 -- Aug 2025}}  

\bullets  

  \item Built \textbf{semantic product search} using Amazon Titan text embeddings and Elasticsearch vector search to enable \textbf{natural language product discovery}, supporting improved user experience for Java-based applications.  

  \item Engineered a \textbf{CDC pipeline} leveraging Flink, Kafka, Debezium, PostgreSQL, RocksDB, and Elasticsearch that reduced synchronization latency from 10–15 minutes to approximately 500 ms while ensuring \textbf{exactly-once processing semantics}, demonstrating ability to deliver high-quality data engineering solutions.  

  \item Deployed and managed Istio/Envoy service mesh across microservices for resilient \textbf{production-ready cloud infrastructure} and integrated distributed tracing with New Relic to enable proactive system monitoring.  

\finishbullets  

\entry  
  {\textbf{Software Engineer I} $|$ \textbf{Infra.Market, Chennai}}  
  {\textbf{Jul 2023 -- Mar 2024}}  

\bullets  

  \item Owned \textbf{Search and Location microservices}, building scalable gRPC APIs handling 500K+ daily requests, contributing to \textbf{real-time product discovery and geolocation features} with emphasis on quality, reliability, and performance.  

  \item Identified and fixed PostgreSQL connection leaks using Gorm ORM, reducing connection pool exhaustion by 85\%, validated under load testing at 10K RPS to ensure \textbf{stable service performance} in deployment environments.  

  \item Developed efficient Airflow pipelines for \textbf{BigQuery-to-Elasticsearch data migration}, optimizing query performance and automating the monitoring of stale search indices to maintain data freshness aligned with business needs.  

\finishbullets  

%-----------PROJECTS-----------  
\section{Projects}  

\entry  
  {GPU-Optimized Personalized Product Search}  
  {Aug 2026 -- Oct 2026}  

\textit{Python, PyTorch, BGE, NVIDIA Triton, Go, gRPC}  

\bullets  

  \item Developed two-tower retrieval with \texttt{BAAI/bge-small-en-v1.5} query/product embeddings and a trainable user tower that leverages interaction history to provide \textbf{personalized product recommendations} with low latency.  

  \item Built an inference pipeline combining precomputed product embeddings, \textbf{approximate nearest-neighbor retrieval}, and a Go/gRPC gateway integrated with NVIDIA Triton for GPU-accelerated query encoding serving.  

  \item Designed scalability benchmarks for steady and bursty traffic measuring latency, throughput, GPU utilization, and cost; evaluated retrieval performance metrics such as Recall@K and NDCG@K with overload recovery testing.  

\finishbullets  

\entry  
  {Attention-Based Hallucination Detection in NMT}  
  {Apr 2026 -- May 2026}  

\textit{Python, PyTorch, mBART, Transformers, scikit-learn}  

\bullets  

  \item Built English-German hallucination detection leveraging mBART attention features from WMT22/WMT23 datasets; deployed multi-layer attention-based classifiers achieving F1 of 0.712 and ROC-AUC of 0.695.  

  \item Extracted diverse attention metrics (entropy, max attention mass, variance) across decoder layers and heads to improve \textbf{machine translation quality estimation} and enable hallucination-aware reranking approaches.  

  \item Demonstrated AI/ML application in \textbf{NLP model interpretability} and quality control, relevant to real-world AI software solutions.  

\finishbullets  

\entry  
  {Prompt/LoRA for Self-Supervised Learning}  
  {Ongoing}  

\textit{Python, PyTorch, Vision Transformers}  

\bullets  

  \item Researching domain-specific visual prompting and LoRA application in self-supervised ViT frameworks to enhance \textbf{transfer learning and adaptability} for computer vision models under domain shifts.  

  \item Evaluating parameter-efficient fine-tuning strategies and inference performance to improve training and deployment efficiency for vision-based AI.  

\finishbullets  

\end{document}